%%%%%%%%%%%%%%%%%%%%%---Chapter3 Formatting---%%%%%%%%%%%%%%%%%%%%%%%%%%%%%%%%%%
\documentclass[Thesis.tex]{subfiles}
\begin{document}
\newpage
\setcounter{chapter}{2} %%Chapter Numbering

\titleformat{\chapter}[display]
{\sffamily\fontsize{27}{27}\bfseries\filleft}{\thechapter}{0pt}{{#1}}

\thispagestyle{empty}
%%%%%%%%%%%%%%%%%%%%%%%%%%%%%%%%%%%%%%%%%%%%%%%%%%%%%%%%%%%%%%%%%%%%%%%%%%%%

\chapter{Methodology}\label{chap3:methodology}

This chapter formalizes the benchmark. Section~\ref{sec:ch3_formulation} states the
learning problem, Sections~\ref{sec:ch3_noise}--\ref{sec:ch3_dataset} specify the
data-generating process, Section~\ref{sec:ch3_protocol} defines the evaluation
protocol and explains why a particular class of data split is invalid here, and
Sections~\ref{sec:ch3_analytic}--\ref{sec:ch3_hybrid} define the nine model arms.
Section~\ref{sec:ch3_metrics} defines the metrics and the cost accounting. The
methodology for the syndrome-decoding track is different in almost every respect and
is therefore stated separately in Chapter~\ref{chap6:decoding}.

\section{Problem formulation}\label{sec:ch3_formulation}

Let $C$ be a quantum circuit drawn from some family, and $O$ a Pauli observable
measured on its output. Write $\mu_{\text{ideal}}(C, O)$ for the noise-free
expectation value and $\mu_\lambda(C, O)$ for the expectation value measured on a
device whose noise has been scaled by a factor $\lambda \geq 1$, so that
$\lambda = 1$ is the bare device. A mitigation method is a function
\begin{equation}
    \widehat{\mu} \;=\; f\!\left( \mu_{\lambda_1}, \ldots, \mu_{\lambda_K},\; s(C, O) \right),
    \label{eq:ch3_estimator}
\end{equation}
where $s(C, O)$ is any classically available description of the circuit and
observable, and the quality of the method is the closeness of $\widehat{\mu}$ to
$\mu_{\text{ideal}}$ over the family. Analytic methods fix $f$ in advance;
learned methods fit $f$ on pairs
$\left( \mu_{\lambda_1:\lambda_K}, \mu_{\text{ideal}} \right)$ obtained on reference
circuits. This thesis uses $K = 3$ with
$\lambda \in \{1, 2, 3\}$, which is the standard choice for three-point
extrapolation and fixes the quantum cost of every arm that uses all three inputs at
three circuit evaluations per estimate.

The unit of analysis is the pair $(C, O)$, not the circuit. Each circuit is measured
for five observables, so each circuit contributes five rows to the dataset. This
distinction becomes important in Section~\ref{sec:ch3_protocol}.

\section{Noise regime I: idealized Markovian model}\label{sec:ch3_noise}

The benchmark is run under two noise regimes. This section defines the first, an
idealized gate-local Markovian model with homogeneous rates and exact noise
scaling; Section~\ref{sec:ch3_realistic} defines the second, a device-realistic
model with inhomogeneous rates, coherent errors, crosstalk and non-Markovian
dynamics. Reporting both is deliberate: a conclusion that holds in only one of them
is a conclusion about the noise model, and the reader needs to be able to tell which
is which. The idealized regime is presented first because it is the regime in which
essentially all published machine-learning mitigation results are obtained, so it is
the setting in which this thesis can be compared to the literature.

Noise is applied as a gate-level Kraus channel model. After each single-qubit gate a
depolarizing channel acts on the gate's qubit,
\begin{equation}
    \mathcal{E}^{\text{dep}}_{p}(\rho) \;=\; (1 - p)\,\rho \;+\; \frac{p}{3}
    \left( X \rho X + Y \rho Y + Z \rho Z \right),
    \label{eq:ch3_depol}
\end{equation}
with probability $p_1$~\cite{nielsen2012book}. After each two-qubit gate the same channel acts
independently on both qubits with probability $p_2$. At the end of every circuit
layer an amplitude-damping channel
\begin{equation}
    \mathcal{E}^{\text{AD}}_{\gamma}(\rho) \;=\; K_0 \rho K_0^{\dagger} + K_1 \rho K_1^{\dagger},
    \qquad
    K_0 = \begin{pmatrix} 1 & 0 \\ 0 & \sqrt{1 - \gamma} \end{pmatrix},
    \quad
    K_1 = \begin{pmatrix} 0 & \sqrt{\gamma} \\ 0 & 0 \end{pmatrix},
    \label{eq:ch3_ad}
\end{equation}
acts on every qubit with rate $\gamma$, representing $T_1$ relaxation during the
layer. The base strengths are $p_1 = 0.002$, $p_2 = 0.02$ and $\gamma = 0.004$; the
two-qubit rate is an order of magnitude above the single-qubit rate, which reflects
the ratio typical of superconducting hardware.

Noise scaling is performed in the noise \emph{parameters} rather than by circuit
folding: at scale $\lambda$ the channel strengths become
$(\lambda p_1, \lambda p_2, \lambda \gamma)$, clipped to remain valid probabilities.
Parameter scaling is the appropriate choice for a simulation study because it is
exact -- the simulator is asked for a device that is precisely $\lambda$ times
noisier -- whereas unitary folding and pulse stretching are hardware techniques that
approximate this idealization~\cite{temme2017error,giurgicatiron2020digital,he2020zero}.
In this thesis ``parameter scaling'' always means this multiplication of the
\emph{channel} parameters. It is not the parameter-noise scaling of
Mitiq~\cite{larose2022mitiq} (Section~\ref{sec:ch2_noise_amplification}), which
injects additional random noise into gate rotation angles. The consequence for interpretation is stated in
Section~\ref{sec:ch3_validity}: results here describe the behaviour of mitigation
methods given ideal noise scaling, and a hardware study would additionally inherit
the error of the scaling technique itself.

\section[Noise regime II: device-realistic non-Markovian model]{Noise regime II: device-realistic model}
\label{sec:ch3_realistic}

The idealized model of Section~\ref{sec:ch3_noise} is gate-local, Markovian,
homogeneous across qubits and free of coherent error. Real superconducting
processors violate all four assumptions. This section defines a second regime that
relaxes each of them, so that the benchmark can be run twice and the robustness of
every conclusion to the noise model can be read off directly.

\subsection{Inhomogeneity, coherent error and crosstalk}

Three modifications make the incoherent part of the model device-like.

\emph{Per-qubit inhomogeneity.} Each rate is drawn once per device realization as
\[
  r_q = \bar{r} \cdot \max\bigl(0.3,\ \min(2.0,\ 1 + 0.35\,\xi_q)\bigr),
  \qquad \xi_q \sim \mathcal{N}(0,1),
\]
giving a spread of roughly a factor of three between
the best and worst qubit, as calibration data typically show. Base rates are
$\bar{p}_1 = 0.0015$, $\bar{p}_2 = 0.012$, $\bar{\gamma} = 0.003$ and a separate
pure-dephasing rate $\bar{\gamma}_\phi = 0.004$, the last representing $T_2$ decay
independent of $T_1$ and absent from regime~I. The realized per-qubit values are
recorded in the dataset metadata so the device is reproducible.

\emph{Coherent over-rotation.} Every rotation angle is miscalibrated by a fixed
fractional error, $\theta \mapsto \theta(1 + \epsilon_q)$ with
$\epsilon_q \sim \mathcal{N}(0, 0.02^2)$ drawn once per qubit. This is a systematic,
unital, non-stochastic error: it does not decohere the state and cannot be described
by a Pauli channel, so it is invisible to the noise-scaling assumption that
underlies extrapolation. Coherent errors are the component of hardware noise that
Pauli-twirling is designed to remove and that error mitigation handles worst.

\emph{Always-on ZZ crosstalk.} Between each neighbouring pair on the ring an
unwanted $\exp(-i J_q Z_a Z_b)$ term acts every layer, with $J_q$ inhomogeneous
around $\bar{J} = 0.015$. This is the static coupling that persists whether or not a
gate is being applied, and it correlates errors between qubits -- so the noise is no
longer gate-local.

\subsection{Non-Markovian dynamics from a coupled defect}

The Markov approximation assumes the environment forgets faster than the system
evolves, so that the reduced dynamics is a composition of independent channels. The
canonical microscopic violation of this in superconducting qubits is a two-level-system
(TLS) defect: a long-lived environmental mode that couples coherently to a qubit,
absorbs excitation and returns it~\cite{mueller2019tls,klimov2018fluctuations}.

This is modelled with a collision construction. One additional qubit, wire $n$,
represents the defect. It is initialized in $\ket{0}$ once at the start of the
circuit and \textbf{never reset}. At the end of every layer it exchanges energy with
qubit $0$ through an XY interaction
\begin{equation}
    U_{\text{TLS}}(g) \;=\; \exp\!\left[ -i g \left( X_0 X_{n} + Y_0 Y_{n} \right) \right],
    \label{eq:ch3_tls}
\end{equation}
after which the defect itself relaxes weakly at rate $\gamma_{\text{TLS}} = 0.01$.
The production coupling is $g = 0.15$~rad per layer. Because the defect retains
phase information between layers, the map on the system alone is no longer a
composition of independent channels: excitation that leaves qubit $0$ at one layer
can return at a later one. Tracing out wire $n$ at the end of the circuit gives the
reduced system state from which all observables are computed, exactly as a
laboratory measurement would.

\subsection{Quantifying the non-Markovianity}\label{sec:ch3_blp}

Non-Markovianity is asserted in much of the mitigation literature and measured in
little of it, so it is measured here. The criterion is that of Breuer, Laine and
Piilo~\cite{breuer2009measure}: for a Markovian (more precisely, CP-divisible) map
the trace distance
\begin{equation}
    D(\rho_1, \rho_2) \;=\; \tfrac{1}{2} \left\| \rho_1 - \rho_2 \right\|_1
    \label{eq:ch3_tracedist}
\end{equation}
between the evolved images of any two initial states is non-increasing, because a
channel cannot increase distinguishability. Any revival is therefore information
returning from the environment, and the total revival
\begin{equation}
    \mathcal{N} \;=\; \max_{\rho_1, \rho_2} \sum_{l} \max\!\left\{ 0,\;
    D\!\left(\rho_1^{(l+1)}, \rho_2^{(l+1)}\right)
    - D\!\left(\rho_1^{(l)}, \rho_2^{(l)}\right) \right\}
    \label{eq:ch3_blp}
\end{equation}
is zero for a Markovian map and positive otherwise, where $l$ indexes layers and the
maximization runs over a set of initial state pairs. Alternative measures based on
CP-divisibility or channel capacities exist and need not agree in general; the
trace-distance measure is used here because it has a direct operational reading as
information back-flow~\cite{rivas2014nonmarkov}.

One methodological point matters here and is easy to get wrong. Non-Markovianity is
a property of the \emph{noise map}, not of the circuit, so $\mathcal{N}$ must be
measured on an idling register -- noise channels, crosstalk and defect collisions
only, with no gates. Running the probe through a scrambling circuit instead measures
unitary mixing: the probe is the reduced state of one qubit, and in an entangling
circuit the other register qubits act as an environment for it, so its trace distance
revives periodically even under strictly Markovian noise (for the whole register the
distance would be invariant under the circuit unitaries), and the estimator then reports a
spurious $\mathcal{N} > 0$ for a model that is Markovian by construction. The
measurement in Section~\ref{sec:ch5_nonmarkov} therefore follows the protocol of a
device memory experiment, and both memoryless controls defined below return $\mathcal{N} = 0$
exactly, which is the check that the estimator is sound.

\subsection{Hardware-realizable noise scaling: unitary folding}
\label{sec:ch3_folding}

Regime~I scales noise by multiplying channel strengths, which no device can do.
Regime~II therefore also provides the scaling method practitioners actually
use~\cite{giurgicatiron2020digital,larose2022mitiq}: unitary folding, which
amplifies noise by lengthening the circuit while preserving its ideal action. At the
layer level, a folded layer $L$ is replaced by
\begin{equation}
    L \;\longmapsto\; L \, L^{\dagger} L,
    \label{eq:ch3_fold}
\end{equation}
which is logically the identity composed with $L$ but passes three times as much
noise. Folding every layer of a depth-$d$ circuit gives $\lambda = 3$; folding none
gives $\lambda = 1$; intermediate $\lambda$ is obtained by folding a fraction
$(\lambda - 1)/2$ of the layers, so $\lambda = 2$ folds half of them. The inverse
layer is itself noisy, as it would be on hardware, and the always-on crosstalk and
defect collisions are not inverted -- they act during the folded segments too.

Every row of the realistic dataset therefore carries two independent feature
triples: $\lambda \in \{1,2,3\}$ by parameter scaling and by unitary folding.
Section~\ref{sec:ch5_scaling} compares them, which answers a question the idealized
regime cannot: whether the learned advantage is an artifact of an unrealizable
scaling method.

\subsection{Calibration drift and the paired-control design}\label{sec:ch3_drift}

A learned mitigator is trained against one calibration of a device and used against
a later one. To test that directly, a second calibration of the same device is
generated by multiplying every per-qubit incoherent rate (single- and two-qubit
depolarizing, amplitude damping, dephasing) and every crosstalk strength by an
independent factor $1 + 0.3\,\xi$, $\xi \sim \mathcal{N}(0,1)$ clipped to
$[0.2, 3]$, and by adding a further Gaussian perturbation of standard deviation
$0.3\,\epsilon$ to each coherent over-rotation. The defect parameters are left
unchanged. This represents a recalibration some hours later. Drift of this magnitude in relaxation
times is documented on real devices~\cite{klimov2018fluctuations}.

Four datasets are produced, all sharing identical circuit parameters and identical
ideal expectation values so that every comparison is paired and the only thing that
changes is the noise:
\begin{itemize}
    \item \textbf{A} -- the production realistic model (calibration A);
    \item \textbf{B} -- the same device after drift (calibration B);
        \item \textbf{M} -- a \emph{defect-removed} control identical to A in
    inhomogeneity, coherent error and crosstalk, with the defect decoupled
    ($g = 0$);
    \item \textbf{R} -- a \emph{reset} control identical to A in every respect,
    including the defect coupling, except that the defect is returned to $\ket{0}$
    at the end of every layer. It inflicts the same damage on qubit~$0$ within each
    layer but cannot carry an excitation to the next.
\end{itemize}
Both controls are Markovian, but they are not equivalent. Removing the defect
removes its energy loss as well as its memory, so A against M measures the two
together; A against R holds the per-layer loss fixed and isolates memory. The
distinction turns out to matter (Section~\ref{sec:ch5_markovianity}). Comparing a
model trained on A and tested on B isolates the effect of drift. Under parameter
scaling the defect coupling and decay are multiplied by $\lambda$ together with
every other channel, in A and R alike. Neither
comparison is available in the idealized regime.

\section{Circuit families and observables}\label{sec:ch3_circuits}

Two circuit families are used, chosen because they represent the two workloads for
which mitigation is usually deployed.

The \textbf{hardware-efficient ansatz} (HEA) family applies, in each layer, an
$R_Y(\theta)$ and an $R_Z(\theta)$ rotation to every qubit followed by a ring of
CNOT gates. Rotation angles are drawn independently and uniformly from
$[-\pi, \pi]$. This is the circuit shape used by variational algorithms~\cite{cerezo2021vqa}.

The \textbf{Trotterized transverse-field Ising} (TFIM) family applies, in each layer,
an $R_X(\theta)$ rotation to every qubit followed by a ring of $R_{ZZ}(\theta)$
couplings, which is one Trotter step of evolution under a transverse-field Ising
Hamiltonian. This is the circuit shape used by quantum simulation workloads.

Both families are generated at depths $d \in \{2, 3, 4, 5, 6\}$ layers. Five Pauli
observables are measured on every circuit: $Z_0$, $Z_0 Z_1$, $X_0 X_1$,
$Z_1 Z_2 Z_3$, and the full-weight string $Z_0 Z_1 \cdots Z_{n-1}$. All five have
eigenvalues $\pm 1$, which makes the finite-sampling model of
Section~\ref{sec:ch3_dataset} exact, and their Pauli weights span $1$ to $n$, which
allows the dependence of mitigation quality on observable locality to be measured.

\section{Dataset construction}\label{sec:ch3_dataset}

For each circuit the density matrix is computed by exact simulation at
$\lambda \in \{0, 1, 2, 3\}$, where $\lambda = 0$ denotes the noise-free circuit and
supplies the regression target. Expectation values follow as
$\mathrm{Tr}[O\rho]$. Because the full density matrix is available, all five
observables are obtained from a single simulation per $(\text{circuit}, \lambda)$
pair.

Finite sampling is then applied on top of the exact values. For a $\pm 1$-eigenvalue
observable with exact expectation $\mu$, the probability of the $+1$ outcome is
$q = (1 + \mu)/2$, so an $N$-shot estimate is
\begin{equation}
    \widehat{\mu}_N \;=\; \frac{2k}{N} - 1, \qquad k \sim \mathrm{Binomial}(N, q),
    \label{eq:ch3_shots}
\end{equation}
which is exactly the sampling distribution of the measured estimator. The default
budget is $N = 4096$ shots per expectation value. Both the exact and the sampled
values are stored, which allows the shot-budget sweep of
Section~\ref{sec:ch5_shots} to be performed by re-sampling
Eq.~\eqref{eq:ch3_shots} without re-simulating any circuit.

The primary dataset comprises $900$ circuits ($450$ per family) at $n = 5$ qubits,
giving $4500$ rows. Five secondary datasets, each of $500$ circuits, support the
sweeps: three at $n = 5$ with the noise strengths scaled by $0.5\times$, $2\times$
and $4\times$, and two at $n = 4$ and $n = 6$ qubits with the base noise.

\section{Feature construction}\label{sec:ch3_features}

Every learned arm receives the same six features: the three noise-scaled
expectation values $\mu_{\lambda=1}$, $\mu_{\lambda=2}$, $\mu_{\lambda=3}$, and
three classically available structural descriptors -- the circuit depth, the
two-qubit gate count, and the Pauli weight of the observable. The first three are
exactly the inputs an analytic extrapolation uses, which makes the comparison
between learned and analytic arms a comparison of estimators over identical
measurement data rather than a comparison of information. The structural features
cost nothing quantum: they are read off the circuit description.

Features are standardized to zero mean and unit variance using statistics computed
on the training fold only, and the same transform is applied to the test fold. The
feature ablation of Section~\ref{sec:ch5_ablation} varies this set to determine how
much of the learned advantage comes from the noise-scaled measurements and how much
from circuit structure.

\section{Evaluation protocol}\label{sec:ch3_protocol}

\subsection{Why the row-wise split is invalid}

The five rows contributed by one circuit share that circuit's parameters, its depth,
and -- critically -- its single noise realization. A split that assigns rows
independently to training and test folds will therefore place some observables of a
given circuit in training and others in test, and a model can exploit the shared
structure to predict the held-out rows. The resulting estimate of generalization
error is optimistic and does not describe the deployment situation, in which an
entirely unseen circuit arrives. Every split in this thesis is therefore
\emph{grouped by circuit}: all five rows of a circuit are assigned to the same fold,
and this grouping is also respected inside every cross-validation loop used for
hyperparameter selection.

\subsection{The three splits}

\textbf{Split A, held-out random circuits.} A $70/30$ grouped split of the primary
dataset. This measures in-distribution performance and answers the question, can a
model learn this device's error map from reference circuits of the same kind?

\textbf{Split B, depth extrapolation.} Train on all circuits of depth
$d \leq 4$, test on all circuits of depth $d \geq 5$. This is the deployment
question: mitigation models are trained on circuits shallow enough to reference and
applied to deeper ones. Note that this split is defined by the design, not sampled,
so it has no randomness.

\textbf{Split C, cross-family transfer.} Train on one circuit family and test on the
other, in both directions. This measures whether a learned error map transfers
across workload types on the same device.

\section{Analytic baselines}\label{sec:ch3_analytic}

Four arms require no training.

The \textbf{unmitigated} estimator is $\widehat{\mu} = \mu_{\lambda = 1}$: the raw
measurement. It is the reference against which error reduction is quoted, and its
error defines what a method must beat to be worth deploying at all.

\textbf{Linear zero-noise extrapolation} fits a degree-one polynomial to
$\{(\lambda_k, \mu_{\lambda_k})\}$ by least squares and evaluates it at
$\lambda = 0$. \textbf{Richardson extrapolation}~\cite{temme2017error} of degree
two fits a quadratic,
which with $K = 3$ points is exactly determined. In general, for a polynomial of
degree $m$ fitted to $m+1$ noise levels, the estimator is the linear combination
\begin{equation}
    \widehat{\mu}_{\text{ZNE}} \;=\; \sum_{k=0}^{m} c_k\, \mu_{\lambda_k},
    \qquad
    c_k \;=\; \prod_{j \neq k} \frac{\lambda_j}{\lambda_j - \lambda_k},
    \label{eq:ch3_richardson}
\end{equation}
whose coefficients satisfy $\sum_k c_k = 1$ and alternate in sign. The alternation
is the source of the method's characteristic weakness: the coefficients have
magnitude greater than one, so the estimator amplifies the statistical noise of the
inputs even as it cancels the systematic bias. Appendix~\ref{AppendixA} derives
Eq.~\eqref{eq:ch3_richardson} and its variance inflation.

\textbf{Exponential extrapolation} fits $\mu(\lambda) = a e^{-b\lambda} + c$ by
non-linear least squares, motivated by the observation that under global
depolarizing noise the expectation value decays exponentially in circuit
volume~\cite{endo2018practical,cai2021multiexponential}. With
three data points and three free parameters the fit is exactly determined, and is
included here specifically to show what that costs.

\section{Classical machine-learning arms}\label{sec:ch3_classical}

Three classical regressors are trained on the six features of
Section~\ref{sec:ch3_features}.

\textbf{Ridge regression} with the regularization strength selected by
cross-validation over a logarithmic grid. It is the minimal learned model and
establishes how much of the problem is linear.

\textbf{Gradient-boosted regression trees}, in the histogram-based formulation, with
the number of leaves and the learning rate selected by grouped cross-validation. Its
piecewise-constant, axis-aligned inductive bias means it extrapolates by holding the
boundary value constant, which -- as Chapter~\ref{chap5:results} shows -- turns out
to matter a great deal.

\textbf{A feed-forward neural network} with two hidden layers of $64$ units,
trained to convergence with early stopping on an internal validation split.

\section{Hybrid quantum--classical arms}\label{sec:ch3_hybrid}

\subsection{Quantum fidelity-kernel ridge regression}

The first hybrid arm embeds each standardized feature vector $x$ into a quantum state
$\ket{\phi(x)}$ using a two-repetition IQP-style embedding on six qubits, and defines
the kernel as the squared overlap of two such
states~\cite{havlicek2019supervised,schuld2021kernel},
\begin{equation}
    k(x, x') \;=\; \left| \braket{\phi(x) | \phi(x')} \right|^2 .
    \label{eq:ch3_fidelity}
\end{equation}
This is the fidelity kernel. Ridge regression is then performed in the induced
reproducing-kernel Hilbert space, which requires only the Gram matrix and never the
explicit feature map. Because the embedding is simulated exactly in this work, the
Gram matrix is computed from the stored state vectors as a single matrix product,
which is what makes the arm tractable at this dataset size.

A \textbf{bandwidth} hyperparameter $\beta$ rescales the input before embedding,
$x \mapsto \beta x$. Its role is not cosmetic. As the embedded states spread over a
large fraction of Hilbert space, off-diagonal fidelities collapse towards zero, the
Gram matrix approaches the identity, and the regressor degenerates to memorizing the
training set with no ability to generalize. Bandwidth controls the kernel's spectrum
and hence its inductive bias, so it is selected here by grouped cross-validation
over $\beta \in \{0.05, 0.15, 0.4, 1.0\}$ jointly with the ridge
penalty~\cite{shaydulin2022importance,canatar2023bandwidth}.
Section~\ref{sec:ch5_bandwidth} reports what happens when it is left at $\beta = 1$.

\subsection{Data-reuploading variational circuit}

The second hybrid arm is a data-reuploading variational quantum
circuit~\cite{perezsalinas2020reuploading} on four
qubits with four layers. In each layer $\ell$ and on each qubit $q$ the circuit
applies
\begin{equation}
    R_Y\!\left( w_{\ell q} \cdot x + b_{\ell q} \right) R_Z\!\left( b_{\ell q} \right),
    \label{eq:ch3_reupload}
\end{equation}
followed by a ring of CNOT gates; the data vector $x$ therefore re-enters the circuit
at every layer, which is what gives the model its expressivity~\cite{schuld2021effect}. The prediction is an
affine function $a \langle Z_0 \rangle + c$ of the expectation value of $Z_0$, with
$a$, $c$, $\{w_{\ell q}\}$ and $\{b_{\ell q}\}$ all trainable. Training uses the Adam
optimizer on the mean squared error, with gradients obtained by automatic
differentiation through the simulator.

\subsection{Bounded-encoding variant}\label{sec:ch3_bounded}

Eq.~\eqref{eq:ch3_reupload} has a property that is easy to overlook: the map from
feature to rotation angle is \emph{periodic}. A feature value outside the range seen
during training produces an angle that wraps around the circle, so the model's output
does not merely become inaccurate -- it becomes uncontrolled, and can take any value
in its range. This suggests a specific mechanism for out-of-distribution failure and
therefore a specific test. The bounded variant replaces the encoding with
\begin{equation}
    R_Y\!\left( \frac{\pi}{2} \tanh\!\left( w_{\ell q} \cdot x + b_{\ell q} \right) \right),
    \label{eq:ch3_bounded}
\end{equation}
which confines the angle to $[-\pi/2, \pi/2]$, the region in which $\sin$ is
monotone. Out-of-range inputs now saturate rather than wrap. The two variants are
identical in every other respect -- same qubit count, layer count, parameter count,
optimizer and schedule -- so the comparison isolates the effect of encoding
periodicity. This is the experiment of Section~\ref{sec:ch5_bounded}.

\section{Metrics and cost accounting}\label{sec:ch3_metrics}

The primary metric is the root-mean-square error of the mitigated estimator against
the ideal expectation value over the test fold,
\begin{equation}
    \mathrm{RMSE} \;=\;
    \sqrt{ \frac{1}{M} \sum_{i=1}^{M}
    \left( \widehat{\mu}_i - \mu_{\text{ideal}, i} \right)^2 } .
    \label{eq:ch3_rmse}
\end{equation}
Mean absolute error and the coefficient of determination $R^2$ are reported
alongside. Uncertainty on the RMSE is obtained by bootstrap resampling \emph{over
test circuits} rather than over rows, so that the resampling unit matches the
grouping unit of the split; $95\%$ intervals are reported as percentiles of the
bootstrap distribution. Error reduction is quoted relative to the unmitigated
estimator on the same fold,
$100 \times (1 - \mathrm{RMSE}_{\text{method}} / \mathrm{RMSE}_{\text{unmitigated}})$,
so a negative value means the method is worse than not mitigating at all.

Whether two methods differ is decided by a \emph{paired} bootstrap. Two methods
evaluated on the same test circuits make correlated errors, so whether their marginal
intervals overlap is not a test of a difference between them. Instead the test
circuits are resampled with replacement, all rows of a circuit together, and on each
of $1000$ resamples the difference $\mathrm{RMSE}_a - \mathrm{RMSE}_b$ is recomputed;
the difference is called \emph{resolved} when the $95\%$ percentile interval of this
distribution excludes zero. The same construction compares one method across two
noise conditions, since the datasets share circuits and ideal values. Because a
single split can favour one method by chance, every headline comparison is also
repeated over ten independent circuit-grouped splits, and the variational arms over
ten initialisation seeds.

Cost is reported on three axes. \textbf{Circuit evaluations per prediction} is the
quantum cost and the one that matters operationally: one for the unmitigated
estimator and three for every method consuming the three noise-scaled inputs. The
variational arms add one circuit of their own per prediction. The fidelity kernel
adds one inversion-test circuit per \emph{training point} per prediction, because
the prediction needs the kernel between the test point and every training point,
plus $M(M+1)/2$ circuits once to build the training Gram matrix; on hardware each of
these is estimated from shots (Appendix~\ref{sec:appA_kernelcost}). \textbf{Trainable parameters} measures model capacity, with the caveat that
kernel ridge regression has one dual coefficient per training point, so its
prediction cost grows with the training set rather than being fixed.
\textbf{Training time} is reported for completeness but is a classical,
one-off cost and should not be weighed against the per-prediction quantum cost.

\section{Threats to validity}\label{sec:ch3_validity}

The limitations of the design are stated here so that the results of
Chapter~\ref{chap5:results} are not over-read. Because the benchmark is run under
two noise regimes, it is useful to separate the limitations that regime~II removes
from those that remain in both.

\subsection*{Addressed by the realistic regime}

The idealized regime of Section~\ref{sec:ch3_noise} assumes noise that is
gate-local, Markovian, homogeneous across qubits, purely incoherent, and static,
whereas real devices exhibit crosstalk, non-Markovian memory, coherent control error
and calibration drift~\cite{harper2020efficient,nielsen2021gst}. Regime~II
(Section~\ref{sec:ch3_realistic}) relaxes each of these: rates are inhomogeneous,
rotations carry systematic over-rotation, an always-on $ZZ$ term couples neighbours,
a coupled defect supplies verified memory, and a second calibration supplies drift.
Likewise, noise scaling in regime~I is exact by construction, which favours the
extrapolation baselines relative to hardware practice; regime~II additionally
supplies unitary folding, the technique a device can actually execute. Where the two
regimes disagree, the realistic one should be believed, and
Chapter~\ref{chap5:results} reports both so the reader can see where they disagree.

\subsection*{Remaining in both regimes}

First, the data are still simulated. The realistic model is informed by the
qualitative structure of superconducting-device noise but is not fitted to
tomographic or calibration data from a specific processor, so the \emph{magnitudes}
of the errors reported here are not device predictions. What the realistic regime
establishes is the direction and approximate size of the effect that each realism
ingredient has on the ranking of methods.

Second, the defect model supplies one memory mechanism with a single time scale.
Real non-Markovianity is a superposition of many mechanisms -- a spectrum of
defects, $1/f$ flux noise, leakage to non-computational states -- and the measured
$\mathcal{N}$ characterizes the model used here, not a device.

Third, the system sizes are small ($n \leq 6$ system qubits, plus the defect in
regime~II), bounded by the cost of exact density-matrix simulation, and the ideal
reference values require classical simulability. The sweeps of
Sections~\ref{sec:ch5_noise}--\ref{sec:ch5_qubits} probe the direction of the trends
in noise strength and system size, but a claim about the behaviour of these methods
at utility-scale circuit sizes is not something this thesis is in a position to
make.

Fourth, the hybrid arms are trained by backpropagation through the simulator and, in
the idealized regime, evaluated with exact Gram matrices. The shot-resolved study of
Section~\ref{sec:ch5_shotkernel} quantifies what finite sampling costs them, but
hardware training of the variational arm -- parameter-shift gradients under shot
noise -- is not attempted.

\end{document}
